%!TEX TS-program = xelatex
%!TEX encoding = UTF-8 Unicode
% chktex-file 1

% Persian edition of the resume.
% Architecturally identical to english/nedamani-resume-en.tex: same class
% options, same page geometry, same spacing knobs, same section order, same
% modular cv/*.tex split. The language-specific parts are kept to a minimum
% and marked below.

% This template has been downloaded from:
% https://github.com/posquit0/Awesome-CV

% A4 paper size by default, use 'letterpaper' for US letter
\documentclass[11pt, a4paper]{awesome-cv}
% Configure page margins with geometry
\geometry{left=1.4cm, top=.8cm, right=1.4cm, bottom=1.0cm, footskip=.5cm}

% Set false if you don't want to highlight section with awesome color
\setbool{acvSectionColorHighlight}{true}

% If you would like to change the social information separator from a pipe (|) to something else
\renewcommand{\acvHeaderSocialSep}{\enspace\textbar\enspace}

%-------------------------------------------------------------------------------
% TYPOGRAPHY
%-------------------------------------------------------------------------------
% Mirrors english/nedamani-resume-en.tex. microtype must be loaded before
% xepersian (below), since xepersian has to be the last package.
\usepackage{microtype}

% Line-breaking policy.
% A CV is a list of short items, bullets and proper nouns, so a mid-word break
% ("architec- / tures") reads much worse than a slightly looser line. Turning
% hyphenation off is not enough on its own though: TeX would then have no legal
% break point and would let lines run into the margin (overfull \hbox).
% Disabling hyphenation while letting the word spaces absorb the slack
% (\emergencystretch) keeps the text justified *and* inside the text block.
\hyphenpenalty=10000    % no automatic mid-word breaks
\exhyphenpenalty=10000  % never break at an existing hyphen (front-end, AI-assisted, ...)
\tolerance=9999         % accept the natural, stretched line breaks from above
\emergencystretch=1em   % glue available to TeX before a line would overflow

%-------------------------------------------------------------------------------
% RTL, PERSIAN TYPOGRAPHY AND NUMERALS
%-------------------------------------------------------------------------------
% xepersian must be loaded after every other package: it loads bidi, which
% patches the directionality primitives of the whole document.
\usepackage{xepersian}

% Vazirmatn is the single typeface of the document. The Latin font is Vazirmatn
% as well, so a Persian sentence containing TypeScript, React or RESTful API
% shares one type design instead of switching to a visually unrelated family.
%
% Vazirmatn ships no italic file, and the CV styles ask for \itshape (address)
% and \slshape (entry dates, locations). Those switches reach the *Persian*
% font as soon as a Persian glyph appears inside an italicised run, so the
% italics have to be declared here as well as in awesome-cv/fonts.sty --
% otherwise LaTeX silently falls back to upright and warns.
\settextfont[
  Scale=1,
  BoldFont=Vazirmatn-Bold,
  ItalicFont=Vazirmatn-Regular,
  ItalicFeatures={FakeSlant=0.2},
  BoldItalicFont=Vazirmatn-Bold,
  BoldItalicFeatures={FakeSlant=0.2},
  Script=Arabic
]{Vazirmatn}

\setlatintextfont[
  Scale=1,
  BoldFont=Vazirmatn-Bold,
  ItalicFont=Vazirmatn-Regular,
  ItalicFeatures={FakeSlant=0.2},
  BoldItalicFont=Vazirmatn-Bold,
  BoldItalicFeatures={FakeSlant=0.2}
]{Vazirmatn}

% Digits. Vazirmatn carries Latin digits, so keeping the digit font in the same
% family means the phone number, the page number and the Solar Hijri years all
% share one type design. Persian prose still uses Persian digits, which is what
% the dates in cv/*.tex rely on.
\setdigitfont[
  BoldFont=Vazirmatn-Bold,
  ItalicFont=Vazirmatn-Regular,
  ItalicFeatures={FakeSlant=0.2},
  BoldItalicFont=Vazirmatn-Bold,
  BoldItalicFeatures={FakeSlant=0.2}
]{Vazirmatn}

%-------------------------------------------------------------------------------
%	PERSONAL INFORMATION
%	Comment any of the lines below if they are not required
%-------------------------------------------------------------------------------
% Available options: circle|rectangle,edge/noedge,left|right
% `left' puts the photo in the first of the three header minipages, which under
% bidi is the right edge of the page - the correct position for an RTL CV.
\photo{./images/nedamani.jpg}
\name{علیرضا محمدزاده}{ندامانی}
\position{توسعه‌دهنده وب}
\address{تهران، ایران}

% The social registry wraps every display value in \lr{} itself (see
% awesome-cv/social.sty), so pass the raw number here: \mobile{} strips the
% spaces, dashes and parentheses out of the value to build the tel: URI, and an
% \lr{} wrapper inside the argument would end up inside the link target.
\mobile{(+98) 9104844924}
\email{alireza3205@gmail.com}
\homepage{nedamani.ir}
\github{NedaMani}
\linkedin{NedaMani}
\telegram{NedaManiOfficial}

% PDF metadata: some ATS platforms read document properties for indexing
\hypersetup{%
  pdftitle={علیرضا محمدزاده ندامانی - رزومه توسعه‌دهنده فرانت‌اند},
  pdfauthor={علیرضا محمدزاده ندامانی},
  pdfsubject={رزومه - توسعه‌دهنده فرانت‌اند، Next.js و React},
  pdfkeywords={Front-End Developer, Next.js, React, JavaScript, TypeScript, Tailwind CSS, Bootstrap, Redux, Zustand, TanStack Query, Vitest, DaisyUI, HTML, CSS, Sass, Material UI, Git}
}

\begin{document}

% Print the header with above personal information
% Give optional argument to change alignment(C: center, L: left, R: right)
\makecvheader

% Print the footer with 3 arguments(<left>, <center>, <right>)
% Leave any of these blank if they are not needed
% \today is the Solar Hijri "today" under xepersian, e.g. 6 مهر 1405.
\makecvfooter
  {\today}
  {}
  {\thepage}

%	CV/RESUME CONTENT
%	Each section is imported separately, open each file in turn to modify content
% chktex-file 1

%	SECTION TITLE
\cvsection{خلاصه}

%	CONTENT
\begin{cvparagraph}
توسعه‌دهنده وب نتیجه‌محور با علاقه و کنجکاوی زیاد نسبت به حوزه‌های مختلف وب؛ از معماری‌های مدرن فرانت‌اند گرفته تا زیرساخت و سیستم‌های نرم‌افزاری.
تمرکز اصلی من روی توسعه با TypeScript و ساخت ابزارها، پکیج‌ها و محصولات قابل استفاده است و در پروژه‌های متن‌باز نیز مشارکت داشته‌ام. تجربه‌ی کار روی اپلیکیشن‌های وب و دسکتاپ مقیاس‌پذیر و پرکارایی را دارم و از همکاری با تیم‌های چندتخصصی در محیط‌های Agile/Scrum لذت می‌برم. به یادگیری و استفاده از فناوری‌های مختلف علاقه‌مندم و سعی می‌کنم در کنار توجه به جزئیات فنی، تجربه‌ی کاربر و نیاز واقعی محصول را هم در تصمیم‌های مهندسی در نظر بگیرم.

\end{cvparagraph}

% chktex-file 1

\cvsection{مهارت‌ها}

\begin{cvskills}

\cvskill
{برنامه‌نویسی}
{\lr{TypeScript, JavaScript, Python, Rust, SQL, Bash}}

\cvskill
{بک‌اند}
{\lr{Nodejs, Hono, Django, RESTful API}}

\cvskill
{فرانت‌اند}
{\lr{React, Nextjs, TanStack Start, Tauri}, \lr{Tailwind CSS, Bootstrap, Responsive Design, Progressive Web Application, HTML5, CSS3, SASS}}

\cvskill
{پایگاه داده}
{\lr{MySQL, MongoDB}}

\cvskill
{زیرساخت}
{\lr{Docker, Cloudflare Workers, CI/CD (GitHub Actions), Linux Administration, Nx, Git}}

\cvskill
{مهندسی هوش مصنوعی}
{\lr{AI-Assisted Development, Prompt Engineering, OpenCode, Hermes}}

\cvskill
{اصول و روش‌ها}
{\lr{Agile/Scrum, Monorepo Architecture, SOLID, Functional Programming, Code Review, Test-Driven Development}}

\cvskill
{زبان‌ها}
{انگلیسی (حرفه‌ای)، فارسی (زبان مادری)}

\end{cvskills}

% chktex-file 1
\cvsection{سوابق}

\begin{cventries}

\cventry
{فرانت‌اند}
{\href{https://droppgroup.com}{\lr{Dropp Technologies Group}}}
{} % Location: unused, no city is listed for this entry
{شهریور ۱۴۰۵ --\ اکنون}
{
	\begin{cvitems}
		\item {مشارکت در پروژه‌های توسعه نرم‌افزار به‌عنوان توسعه‌دهنده فرانت‌اند.}
    \item {\lr{TypeScript, Nextjs, TanStack Start, React, Vitest, Date Conversion and Localization, Drag and Drop, Storybooks for components, i18n implementation, State Management, TailwindCSS, Shadcn, Responsive Design}}
		\item {\lr{Documentation, Team Collaboration, Agile/Scrum}}‍
	\end{cvitems}
}

\cventry
{پکیج~\lr{NPM}}
{\href{https://www.npmjs.com/package/iran-estehlal}{\lr{Iran Estehlal}}}
{}
{شهریور ۱۴۰۵}
{
  \begin{cvitems}
    \item {یک کتابخانه بدون وابستگی برای تبدیل تاریخ میلادی به هجری شمسی و برعکس، که تبدیل تقویمی قابل‌اعتماد را برای اپلیکیشن‌های \lr{JavaScript} و \lr{TypeScript} فراهم می‌کند.}
    \item {\lr{TypeScript, Vitest, Date Conversion and Localization, Enforced Commit Conventions, Git Hooks, Github Actions}}
  \end{cvitems}
}

\cventry
{اپلیکیشن دسکتاپ}
{\href{https://github.com/wikimediairan/Perseus}{\lr{Perseus}}}
{}
{تیر ۱۴۰۵ --\ اکنون}
{
  \begin{cvitems}
    \item {یک اپلیکیشن چندسکویی برای ترجمه مقالات ویکی‌پدیای انگلیسی به ویکی‌متن فارسی، با ترجمه مبتنی بر هوش مصنوعی و تشخیص هوشمند پیوندها.}
    \item {\lr{TypeScript, Monorepo Architecture, React, Vitest, AI‑Assisted Development, Prompt Engineering, Responsive Design, TailwindCSS, ShadcnUI, Tauri, Rust, Nodejs, Hono, LLM Integration, Nx, Git Hooks, Github Actions}}
  \end{cvitems}
}

\cventry
{افزونه مرورگر}
{\href{https://github.com/NedaMani/RTLeeBOT}{\lr{RTLeeBOT}}}
{}
{تیر ۱۴۰۵}
{
  \begin{cvitems}
    \item {یک افزونه مرورگر که کمبود پشتیبانی صحیح \lr{RTL} و فونت فارسی دلخواه را در پلتفرم‌های گفت‌وگوی هوش مصنوعی از جمله \lr{ChatGPT, Gemini, Claude, Grok, DeepSeek, Perplexity, Copilot, DuckAI} و \lr{NotebookLM} برطرف می‌کند.}
    \item {\lr{Typescript, React, TailwindCSS, Prompt Engineering}}
  \end{cvitems}
}

\cventry
{پکیج~\lr{NPM}}
{\href{https://www.npmjs.com/package/persian-holidays}{\lr{Persian Holidays}}}
{}
{خرداد ۱۴۰۵}
{
  \begin{cvitems}
    \item {یک پکیج سبک که داده‌های ساختاریافته تعطیلات رسمی ایران را برای اپلیکیشن‌هایی که به نمایش یا پردازش تعطیلات نیاز دارند ارائه می‌دهد.}
    \item {\lr{TypeScript, Vitest, Date Conversion and Localization, Enforced Commit Conventions, Git Hooks, Github Actions}}
  \end{cvitems}
}

\cventry
{افزونه~\lr{Obsidian}}
{\href{https://github.com/karfekr/obsidian-persian-calendar}{\lr{Persian Calendar}}}
{}
{آذر ۱۴۰۴ --\ اکنون}
{
  \begin{cvitems}
    \item {یک افزونه تقویم فارسی برای \lr{Obsidian} با بیش از ۱۳ هزار دانلود و بیش از ۱۲۰ ستاره در \lr{GitHub}، که پشتیبانی از تقویم ایرانی و قابلیت‌های مرتبط با تاریخ را در محیط یادداشت‌برداری فراهم می‌کند.}
    \item {\lr{TypeScript, Vitest, SOLID Principles, Date Conversion and Localization, i18n implementation, BEM Methodology, Responsive Design, Git Hooks, Github Actions}}
    \item {\lr{Code Review and Collaboration}}
  \end{cvitems}
}

\cventry
{تحلیل داده}
{\href{https://github.com/NedaMani/youtube_trending_analysis}{YouTube Trending Videos Data Analysis}}
{}
{Jun. 2025}
{
  \begin{cvitems}
    \item {به تحلیل داده‌های مربوط به ویدیوهای ترند یوتیوب پرداختم تا الگوهای مرتبط با محبوبیت، میزان تعامل و عملکرد محتوا را در دسته‌بندی‌های مختلف شناسایی کنم.}
    \item {\lr{Python, Data Analysis, Data Visualization, Pandas, Numpy, Jupyter Notebook}}
  \end{cvitems}
}

\cventry
{فرانت‌اند}
{\href{https://github.com/PhoenixTask/WebApp}{\lr{PhoenixTask}}}
{}
{خرداد ۱۴۰۴ --\ مهر ۱۴۰۴}
{
  \begin{cvitems}
    \item {ساخت یک اپلیکیشن مدیریت کار الهام‌گرفته از \lr{Trello} و پیاده‌سازی مستقل فرانت‌اند آن به‌عنوان یک پروژه شخصی.}
    \item {\lr{TypeScript, Nextjs, React, Date Conversion and Localization, Drag and Drop, i18n implementation, State Management, TailwindCSS, Responsive Design}}
    \item {\lr{Team Collaboration, Agile/Scrum, Intern Mentoring}}
  \end{cvitems}
}

\cventry
{فرانت‌اند}
{\href{https://tourino.app/}{\lr{Tourino}}}
{}
{دی ۱۴۰۳ --\ اردیبهشت ۱۴۰۴}
{
	\begin{cvitems}
    \item {مشارکت در یک اپلیکیشن وب استارتاپی از طریق بهبود \lr{SEO}، کاهش حجم \lr{bundle} و بازآرایی معماری فرانت‌اند برای دستیابی به کارایی و نگهداری‌پذیری بهتر.}
		\item {\lr{TypeScript, Nextjs, React, Date Conversion and Localization, Drag and Drop, State Management, Responsive Design, MaterialUI, SASS, Git Flow}}
    \item {\lr{Team Collaboration, Agile/Scrum}}
  \end{cvitems}
}

\cventry
{فرانت‌اند}
{\href{https://quera.org/}{\lr{Quera}}}
{}
{دی ۱۴۰۲ --\ اردیبهشت ۱۴۰۳}
{
	\begin{cvitems}
    \item {همکاری با تیمی چهار نفره در ساخت یک اپلیکیشن مدیریت کار مشابه \lr{Trello}، با تمرکز بر توسعه فرانت‌اند، حل مسئله و همکاری چابک.}
    \item {\lr{TypeScript, React, Date Conversion and Localization, Drag and Drop, State Management, TailwindCSS}}
    \item {\lr{Team Collaboration, Agile/Scrum}}
  \end{cvitems}
}

\end{cventries}

% chktex-file 1
\cvsection{گواهی‌ها}

\begin{cvhonors}
  \cvhonor
    {\lr{Task-Oriented Bootcamp} در توسعه فرانت‌اند با \lr{React}} % Name
    {کوئرا} % Issuer
    {\href{https://quera.org/certificate/1TyQm8RV/}{گواهی}} % Credential
    {} % Date(s)

  \cvhonor
    {دوره پروژه‌محور توسعه وب فرانت‌اند} % Name
    {کوئرا} % Issuer
    {\href{https://quera.org/certificate/Z4shbXOc/}{گواهی}} % Credential
    {} % Date(s)

  \cvhonor
    {\lr{CS50x}} % Name
    {دانشگاه هاروارد} % Issuer
    {\href{https://certificates.cs50.io/732f87e9-6550-444e-8ad6-99922c3aca9c.pdf}{گواهی}} % Credential
    {} % Date(s)

  \cvhonor
    {دوره تسک‌محور \lr{Linux}} % Name
    {کوئرا} % Issuer
    {\href{https://quera.org/certificate/oBvP6gms/}{گواهی}} % Credential
    {} % Date(s)

  \cvhonor
    {دوره تسک‌محور مبانی برنامه‌نویسی \lr{Python}} % Name
    {کوئرا} % Issuer
    {\href{https://quera.org/certificate/33b4vSHp/}{گواهی}} % Credential
    {} % Date(s)

  \cvhonor
    {دوره تسک‌محور نگارش \lr{Prompt}} % Name
    {کوئرا} % Issuer
    {\href{https://quera.org/certificate/3CfoovY6/}{گواهی}} % Credential
    {} % Date(s)

\end{cvhonors}

% chktex-file 1
%-------------------------------------------------------------------------------
%	SECTION TITLE
%-------------------------------------------------------------------------------
\cvsection{تحصیلات}


%-------------------------------------------------------------------------------
%	CONTENT
%-------------------------------------------------------------------------------
\begin{cventries}

%---------------------------------------------------------
  \cventry
    {کارشناسی ارشد مهندسی نرم‌افزار} % Degree
    {دانشگاه آزاد اسلامی -\ واحد علوم و تحقیقات} % Institution
    {تهران، ایران} % Location
    {۱۴۰۴ --\ اکنون} % Date(s)
    {}

%---------------------------------------------------------
  \cventry
    {کارشناسی مهندسی نرم‌افزار} % Degree
    {دانشگاه ملی مهارت -\ شمسی‌پور} % Institution
    {تهران، ایران} % Location
    {۱۴۰۰ --\ ۱۴۰۴} % Date(s)
    {}

%---------------------------------------------------------
\end{cventries}


\end{document}
